\section{The exact H plumbing}

Let \(M\) be surgery on the standard framed-unknot plumbing link of the weighted tree in Figure~\ref{fig:plumbing}. In this Kirby notation each vertex gives an unknot with the displayed integer framing relative to its zero longitude, and each edge gives a positive Hopf clasp. Precisely, one can start with the left central unknot, add the right central unknot as a meridian of it, and then insert two disjoint meridians of each central component in small neighborhoods away from the previous clasp. Orient the meridians to have linking number \(+1\) with their center. There are no other clasps or twists in this link diagram. In the order of the two centers, the left pair, and the right pair, the matrix is
\[
 Q=\begin{pmatrix}
 0&1&1&1&0&0\\
 1&0&0&0&1&1\\
 1&0&2&0&0&0\\
 1&0&0&-2&0&0\\
 0&1&0&0&2&0\\
 0&1&0&0&0&-2
 \end{pmatrix}.
\]
The two leaf pairs have cancelling Schur-complement corrections, so \(Q\) is congruent over \(\mathbb R\) to the diagonal leaf block plus the hyperbolic plane. Thus \(\det Q=-16\) and \(b_+(Q)=b_-(Q)=3\). In the presentation of \(\operatorname{coker}Q\), eliminate the centers \(a,b\) using \(a=-2u, b=-2w\), for \(u,w\) the \(+2\) leaves, and then eliminate the \(-2\) leaves using the central relations. The two remaining relations are \(4u=4w=0\). In particular, \(H_1(M;\ZZ)\simeq(\ZZ/4\ZZ)^2\). No geometric identification of \(M\) is needed.

\begin{figure}[ht]
\centering
\setlength{\unitlength}{1mm}
\begin{picture}(108,30)(0,0)
 \linethickness{0.7pt}
 \put(32,15){\line(1,0){44}}
 \put(14,24){\line(2,-1){18}}
 \put(14,6){\line(2,1){18}}
 \put(76,15){\line(2,1){18}}
 \put(76,15){\line(2,-1){18}}
 \put(32,15){\circle*{2.5}}
 \put(76,15){\circle*{2.5}}
 \put(14,24){\circle*{2.5}}
 \put(14,6){\circle*{2.5}}
 \put(94,24){\circle*{2.5}}
 \put(94,6){\circle*{2.5}}
 \put(30,18){\makebox(4,5){\(0\)}}
 \put(74,18){\makebox(4,5){\(0\)}}
 \put(3,22){\makebox(8,5){\(+2\)}}
 \put(3,3){\makebox(8,5){\(-2\)}}
 \put(97,22){\makebox(8,5){\(+2\)}}
 \put(97,3){\makebox(8,5){\(-2\)}}
\end{picture}
\caption{The Kirby plumbing diagram defining \(M\). Each edge is a positive Hopf clasp; vertex labels are surgery framings.}
\label{fig:plumbing}
\end{figure}



\begin{theorem}\label{thm:hcase}
For the closed manifold of Figure~\ref{fig:plumbing}, put
\[
 C_r=\sum_{\substack{1\le n<r\\ n\ \mathrm{odd}}}\cot\!\left(\frac{\pi n}{2r}\right).
\]
Then, at every \(r\ge3\),
\begin{equation}\label{eq:exact}
 \tau_r(M)=(-1)^r\frac{C_r}{\sin(\pi/r)},
 \qquad
 Z_r(M)=(-1)^r\sqrt{\frac2r}\,C_r.
\end{equation}
If \(c=\gamma+\log(4/\pi)\), where \(\gamma=\gamma_{\rm E}\) is Euler's constant, then
\begin{align}
 \tau_r(M)&=(-1)^r\frac{r^2}{\pi^2}\bigl[\log r+c+o(1)\bigr],\label{eq:tauasy}\\
 Z_r(M)&=(-1)^r\frac{\sqrt{2r}}{\pi}\bigl[\log r+c+o(1)\bigr].\label{eq:Zasy}
\end{align}
\end{theorem}

\begin{proof}[Proof of Theorem~\ref{thm:hcase}]
The two caps at a trivalent center leave the vector
\[
 u_r(a)=\frac{|(S\Theta^2S)_{a1}|^2}{S_{1a}}
       =\sqrt{2/r}\,\mathbf1_{r-a\ {\rm odd}}
                      \cot\bigl(\pi(r-a)/(2r)\bigr)
\]
by \eqref{eq:cap}. The ordinary signature is zero, so \(Z_r=u_r^{\mathsf T}S u_r\). For \(1\le b<r\) the following odd cotangent transform is exact:
\begin{equation}\label{eq:cot-transform}
 \sum_{n\ {\rm odd}<r}\cot(\pi n/(2r))\sin(\pi nb/r)=r/2.
\end{equation}
For clarity, over the full length-\(N\) grid the transform
\(\sum_{n=1}^{N-1}\cot(\pi n/N)\sin(2\pi jn/N)\) equals \(N-2j\) for \(1\le j<N\). At \(j=1\) the summand is \(1+\cos(2\pi n/N)\), and differences for consecutive \(j\) sum to \(-2\) by root orthogonality. Subtract the even part of the full \(N=2r\) sum and fold about \(r\); the midpoint cotangent is zero. This proves \eqref{eq:cot-transform}. Sine reflection then gives \((Su_r)(b)=(-1)^{b+1}\), which has sign \((-1)^r\) on the support of \(u_r\). The exact formulas follow with \(Z_r=S_{11}\tau_r\).

Subtract \(2r/(\pi n)\) from the odd cotangent summand. The resulting smooth Riemann sum is
\[
 \frac r2\int_0^1\left[\cot(\pi x/2)-\frac{2}{\pi x}\right]dx+O(1)
       =\frac r\pi\log(2/\pi)+O(1).
\]
The subtraction itself sums to \((r/\pi)[\log r+\gamma_{\rm E}+\log2]+O(1)\) on both parities. This gives the stated finite constant and asymptotic estimates.
\end{proof}

